\section{Tokenomics \& Governance}
The Vertex network is powered by the utility and governance token \textbf{VTX}. The token has three core functions: (1) staking for consensus security, (2) payment of network transaction fees, and (3) protocol governance.

\subsection{Token Allocation}
Table 2 outlines the distribution and vesting schedules of the initial VTX token supply.

\begin{table}[H]
\centering
\sffamily
\small
\begin{tabularx}{0.9\textwidth}{Xcr}
\toprule
\textbf{Allocation Category} & \textbf{Percentage} & \textbf{Vesting Schedule} \\
\midrule
Community \& Liquidity Mining & 40.0\% & 48-month linear emission \\
Ecosystem \& Partnerships & 20.0\% & 12-month cliff, 36-month linear \\
Vertex Treasury & 15.0\% & 24-month linear vesting \\
Core Contributors \& Founders & 15.0\% & 12-month cliff, 48-month linear \\
Early Backers \& Seed Round & 10.0\% & 6-month cliff, 18-month linear \\
\bottomrule
\end{tabularx}
\caption{Vertex Token (VTX) distribution and vesting allocations.}
\label{tab:tokenomics}
\end{table}

\subsection{Fee Mechanism}
To maintain long-term token value, Vertex utilizes a dynamic burning mechanism. A fixed portion of the base transaction fee $f_{base}$ is permanently burned, decreasing the circulating supply of VTX as transaction volume increases. The remaining portion of the fee and the user tip are distributed to active validators to incentivize security and network maintenance.
